\begin{lemma}\label{lem:decic333-finite-image}
Put $F=\mathbb Q(\zeta_3)$ and $X=X_1(3,33)$, and let
$h(E,P_3,Q_{33})=(E,P_3,10Q_{33})$.
Then $(h_*-1)\operatorname{Jac}(X)(F)$ is finite.
\end{lemma}
\begin{proof}
Over $F$, the curve $X$ is isomorphic to the base change of
$X_H(99)$, where $H=\langle34\rangle$; see \cite[Section~2]{DS}.
The involution $h$ is induced by the diamond operator
$\langle10\rangle$. The source has genus $81$ and its quotient
by $h$ has genus $41$, so its anti-invariant Jacobian factor
has dimension $40$.

The even nebentypes $\varepsilon$ modulo $99$ satisfying
$\varepsilon(34)=1$ and $\varepsilon(10)=-1$ have conductor
dividing $33$. They have precisely two Galois orbits, of orders
$2$ and $10$. In each case the corresponding cuspidal modular
symbol space with sign $+$ has dimension $8$ over the character
field. Their rational dimensions are $8$ and $4\cdot8$, giving
all $40$ dimensions of the anti-invariant factor.
The exact winding computation gives full Hecke orbit rank
$8$ in both spaces. This is certified by reduction modulo $5$
and $11$, respectively, avoiding the finitely many denominator
primes. Thus every newform contributing to these spaces has
nonzero central $L$-value.

Let $\chi_{-3}$ be the quadratic character of $F$.
The twist-containment argument at the end of the proof of
\cite[Theorem~4.1]{DS}, using Atkin--Li, shows that twisting a
form in $S_2(\Gamma_0(99),\varepsilon)$ by $\chi_{-3}$ stays in
this same space: the conductor of $\varepsilon$ divides $33$
and $99=3\cdot33$. In particular its nebentype and its $h$-sign
are unchanged. The preceding winding computation therefore
also covers every $\chi_{-3}$-twist of an anti-invariant factor.
The nonvanishing theorem for modular abelian varieties used
in \cite[Section~4]{DS} gives rank zero over $\mathbb Q$ for
each factor and its twist, hence rank zero over $F$.
The image of $h_*-1$ is contained in this factor, proving the
claim. The winding vectors and Hecke matrices are retained,
and their full ranks are checked independently in Sage.
\end{proof}

\begin{lemma}\label{lem:decic-relative-formal-F4}
Let $R$ be the integers of the unramified quadratic extension of
$\mathbb Q_2$, with residue field $k=\mathbb F_4$.
Let $X/R$ be a smooth proper curve, let $A=\operatorname{Jac}(X)$,
and let $h$ be an involution acting freely on a neighbourhood of
the cusps. Put $\alpha=h_*-1$ and $\pi:X\to Z=X/\langle h\rangle$.
Suppose that $Z/R$ is smooth near the images of these cusps.
Write a cuspidal divisor $\overline C$ of degree $d$ as
\[
              \overline C=\overline C_0+\pi^*\overline E,
              \qquad \deg\overline E=k_0,
\]
where the multiplicities of $\overline E$ are maximal on each
pair interchanged by $h$. Let $C$ be its cuspidal lift.
Suppose that $\alpha[D-C]$ is torsion for an effective divisor
$D\in X^{(d)}(R)$ with $\overline D=\overline C$.
If restriction of
\[
 U=(\mathcal C^2-1)\alpha^*H^0(X_k,\Omega^1)
\]
to $\overline C$ has rank $d-k_0$, then $D$ contains the fixed
divisor $C_0$.
\end{lemma}
\begin{proof}
We first record the formal-group calculation needed at $2$.
For a torsion point $P$ in the first formal filtration of $A(R)$,
write its coordinates as $t(P)=2u$. The second filtration is
torsion-free, by the formal logarithm, and $[2]P=0$.
If $H$ is the matrix of the quadratic terms of $[2]$, reduction
of $[2](2u)=0$ modulo $8$ gives
\[
                   \overline u=H\overline u^{[2]}.
\]
Squaring and using $\overline u^{[4]}=\overline u$ gives
$\overline u=HH^{[2]}\overline u$.
The linearized Cartier matrix is $H^{[2]}$, so $HH^{[2]}$
is the matrix of $\mathcal C^2$ on cotangents, with the dual
coordinate convention. Thus every form in
$(\mathcal C^2-1)H^0(A_k,\Omega^1)$ annihilates the first
formal tangent of $P$.

Apply this to $P=\alpha[D-C]$.
The first tangent of $D-C$ is therefore annihilated by the
restriction of $U$. Its kernel has dimension $k_0$.
The formal family
\[
                  T=C_0+\pi^*Z^{(k_0)}
\]
through $C$ has dimension $k_0$ and lies in the fibre of
$E\mapsto\alpha[E-C]$.
It is a closed formal immersion into $X^{(d)}$.
Indeed, $\pi$ is \'etale at the cusps. On each pair of local
branches, the two moving monic polynomials are the same, and
maximality of $\overline E$ leaves no fixed factor on at least
one branch. The moving coefficients can therefore be recovered
from that branch. This also proves tangent injectivity.
Its tangent space is consequently the kernel of the restriction
of $U$.

Since $T/R$ is formally smooth, we may choose $C'\in T(R)$
with $C'\equiv D\pmod4$. We have
$\alpha[D-C']=\alpha[D-C]$, because $\alpha\pi^*=0$.
This torsion point lies in the second formal filtration, and
hence is zero.
Finally, the rank of $\alpha^*$ restricted to $\overline C$
is at least $d-k_0$. Choose $d-k_0$ independent cotangent
coordinates of the map $E\mapsto\alpha[E-C]$.
Their common formal zero locus is smooth of dimension $k_0$.
It contains $T$, and the tangent map of $T$ to this locus is an
isomorphism. The formal implicit function theorem therefore
identifies this zero locus with $T$. Since $D$ lies in the full
fibre, $D\in T(R)$, and $C_0\leq D$.
\end{proof}


\begin{proposition}\label{prop:decic-no333}
An elliptic curve over a number field of degree $10$ has no subgroup
isomorphic to $C_3\times C_{33}$.
\end{proposition}
\begin{proof}
Put $F=\mathbb Q(\zeta_3)$ and $X=X_1(3,33)_F$.
The Weil pairing gives $F\subset K$. A proposed level structure
therefore gives an effective noncuspidal $F$-rational divisor $D$
of degree $5$ on $X$. Write $h=\langle10\rangle$ and
$\alpha=h_*-1$. The image $\alpha J_X(F)$ is finite by
Lemma~\ref{lem:decic333-finite-image}.
We use reduction at the prime above $2$, with residue field
$k=\mathbb F_4$.

The Frobenius polynomial and an independent exhaustive enumeration
of elliptic curves give the noncuspidal closed-place counts
\[
                         0,\quad0,\quad0,\quad0,\quad264
\]
in degrees one through five. The closed cusps have signature
$1^{40},5^8$. Thus $\overline D$ is either entirely cuspidal or
one noncuspidal quintic place.

We first exclude the cuspidal case. Put $R=\mathcal O_{F,2}$,
completed at $2$. Cusps lift over unramified extensions of $R$;
for a closed quintic cusp we take the trace of its five geometric
lifts. These traces give an $R$-rational cuspidal lift $C$ of
$\overline D$. Choose an $F$-rational Tate cusp $c_0$.
The class $\alpha[D-5c_0]$ is torsion by finiteness of
$\alpha J_X(F)$, and $[C-5c_0]$ is torsion by
Manin--Drinfeld. Hence $\alpha[D-C]$ is torsion.
This does not require finiteness of $J_X(F)$.

The involution $h$ interchanges the cusps in $20$ rational pairs
and four pairs of closed quintic cusps. Its quotient is \'etale
near every cusp, with smooth local quotient over $R$.
Write
\[
              \overline C=\overline C_0+\pi^*\overline E,
              \qquad k_0=\deg\overline E,
\]
using maximal multiplicities on each pair, where
$\pi:X\to X/\langle h\rangle$.
For a rational cuspidal divisor of degree $5$, we have
$k_0\in\{0,1,2\}$. A single quintic cusp has $k_0=0$.
The space
\[
       U=(\mathcal C^2-1)(h^*-1)H^0(X_k,\Omega^1)
\]
has dimension $36$. Its restriction to every such
$\overline C$ has rank $5-k_0$.
The complete counts for rational cusp divisors are
\[
\begin{array}{c|rrr}
k_0&0&1&2\\ \hline
\text{number}&864008&213600&8400,
\end{array}
\]
and each of the eight quintic cusps has restriction rank $5$.
Lemma~\ref{lem:decic-relative-formal-F4} now implies $C_0\leq D$.
Since $5$ is odd, $C_0$ is nonempty. This contradicts the
noncuspidal support of $D$.

Here the restriction matrices are obtained on the explicit
Kummer presentation over $Y=X_1(33)_k$. If $Q_{33}$ is the
universal point and $11Q_{33}$ is put in the form
$(0,0)$ on
\[
                       y^2+A_1xy+A_3y=x^3,
\]
the additional $3$-level structure is given by
\[
                              r^3=A_1^3-27A_3.
\]
The resulting curve has genus $81$. The exact holomorphic basis,
Cartier matrix and diamond pullback are retained. At all $24$
base cusps we expand the coefficient functions, use the tame
Kummer uniformizer, and obtain the jets at all $48$ closed
source cusps. An independent Sage calculation reconstructs
the function-field equations and these series, and repeats
every restriction matrix and all $1086008+8$ ranks.
The cusp permutation is also checked directly on the canonical
first-jet spaces, which are pairwise distinct.

It remains to exclude noncuspidal reduction. Forgetting the
independent point of order $3$ maps the $264$ quintic places
onto $88$ quintic places on $Y_k$, in twelve diamond orbits.
The Jacobian $J_Y(\mathbb Q)$ is finite by
\cite[Theorem~3.1]{DEHMZ}. For a degree-five divisor $M$ on
$Y_F$, its trace $M+M^\sigma$ is a rational divisor of degree
$10$, where $\sigma$ is the nontrivial automorphism of $F$.
Put
\[
                  A_5=T_5-1-5\langle5\rangle.
\]
Eichler--Shimura and torsion injectivity at $5$ give
$A_5J_Y(\mathbb Q)=0$. At a rational Tate cusp $c_0$,
$T_5c_0=c_0+5\langle5\rangle c_0$ as divisors.
Thus $A_5(M+M^\sigma)$ is principal.

For the twelve possible reduction orbits, the exact
Riemann--Roch tests for this divisor give dimension zero
except at representatives $1$, $6$ and $10$.
At each of these three representatives we also use the three
$3$-isogenies whose kernel is not $\langle11Q_{33}\rangle$.
The image of $Q_{33}$ still has order $33$, so each quotient
gives another divisor $M$ on $Y_F$ to which the same argument
applies. The nine resulting Riemann--Roch dimensions, in order,
are
\[
                       (0,0,1,\;1,0,0,\;1,0,0).
\]
Each representative therefore has a zero-dimensional test,
contradicting principality.

The finite computations retain the residue norms of every
pushforward and trace, including multiplicity two when the
rational trace has one closed place as support.
An independent enumeration checks all $72$ geometric
$5$-kernels and the three selected $3$-kernels, with all $75$
marked V\'elu maps. A second native calculation reconstructs
the $80$ prime references from two-generator ideals, binds
the source and all marked images, and repeats the twelve
selected zero-dimensional Riemann--Roch tests.
\end{proof}
